\documentclass[11pt]{article}
\usepackage[a4paper,margin=2.5cm]{geometry}
\usepackage[T1]{fontenc}
\usepackage{lmodern}
\usepackage[hidelinks]{hyperref}
\setlength{\parindent}{0pt}
\setlength{\parskip}{0.7em}
\pagestyle{empty}

\begin{document}

\begin{flushright}
24 October 2026
\end{flushright}

To the Editor, \emph{Classical and Quantum Gravity}

Dear Editor,

We submit our manuscript \emph{``Rotation slows the instability of Ellis--Bronnikov wormholes''}, by Nikita M.\ Shirokov and Ilya Nachevsky, for consideration as a Paper in \emph{Classical and Quantum Gravity}.

The Ellis--Bronnikov wormhole is unstable: a small perturbation makes its throat collapse to a black hole or expand without bound. Gonz\'alez, Guzm\'an and Sarbach established this in spherical symmetry in this journal (CQG \textbf{26} 015010 and 015011, 2009), and rotation, the oldest candidate for removing it, entered with the rotating scalar-field wormhole of Matos and N\'u\~nez, also in this journal (CQG \textbf{23} 4485, 2006). Perturbation theory to second order in the spin suggests that the instability weakens and may disappear near 30\% of the extremal angular momentum, but it cannot follow the solution beyond that point. The manuscript presents, to our knowledge, the first nonlinear evolutions of rotating wormholes. Its main results are:

\begin{itemize}\setlength{\itemsep}{0.2em}
\item We construct stationary rotating wormholes with a spectral solver up to an equatorial rotation velocity of 99\% of the speed of light, including a family for which no solutions had been published. The solver is validated against published solutions, and checking each solution against field equations the solver does not impose revealed an error in a published field equation, which we correct.
\item Rotation slows the collapse. Its e-folding time rises from 5.4 in the static model to 13.3 at 24\% of the extremal angular momentum, and the collapse rate extrapolates to zero at 32\%, beside the perturbative prediction of 33\%.
\item The instability does not end there. Near 30\% of the extremal angular momentum the unperturbed throat expands instead of collapsing, at a rate that rises with spin.
\item From 58\% of the extremal angular momentum upward we find no growth over the evolved times; at 81\% this holds over about 28 e-folding times of the static instability, at three resolutions.
\item A collapsing rotating throat emits a quadrupolar gravitational-wave burst without any nonspherical perturbation, and the radiated energy grows with spin.
\end{itemize}

The work connects exact stationary solutions, perturbation theory and numerical relativity, and builds on the CQG papers above. We therefore believe it suits the journal's readership.

The solver, the evolution code, the parameter files and the diagnostic data of every run are public (archived at \url{https://doi.org/10.5281/zenodo.23245933}). A preprint is available as arXiv:2610.29941. The manuscript has not been published and is not under consideration elsewhere, and both authors have approved this submission.

Yours sincerely,

Nikita M.\ Shirokov (corresponding author)\\
Moscow Institute of Physics and Technology\\
\href{mailto:shirokov.nm@phystech.edu}{shirokov.nm@phystech.edu}

\end{document}
